% ===== PRÉAMBULE CANONIQUE — à coller VERBATIM en tête de CHAQUE .tex =====
\documentclass[11pt,a4paper]{article}
\usepackage[utf8]{inputenc}\usepackage[T1]{fontenc}\usepackage{lmodern}
\usepackage[french]{babel}
\usepackage{amsmath,amssymb,mathtools}\usepackage{icomma}
\usepackage[version=4]{mhchem}          % \ce{...} : équations & formules
\usepackage{siunitx}\sisetup{locale=FR} % \qty{}{}, \si{}, \ang{}
\usepackage{chemfig}                    % \chemfig{...} : structures développées
\usepackage{xcolor}\usepackage{colortbl}
\usepackage{tikz}\usepackage{pgfplots}\pgfplotsset{compat=1.18}
\usepackage[most]{tcolorbox}\usepackage{enumitem}\usepackage{array,booktabs}
\usepackage[a4paper,margin=2.2cm]{geometry}
\usetikzlibrary{arrows.meta,calc,angles,quotes,positioning,patterns,backgrounds,shapes.geometric}
\definecolor{primary}{RGB}{222,49,99}\definecolor{primarydark}{RGB}{150,22,55}
\definecolor{defcol}{RGB}{0,114,178}\definecolor{methcol}{RGB}{52,92,148}
\definecolor{excol}{RGB}{0,135,90}\definecolor{attncol}{RGB}{200,40,55}
\tcbset{boxbase/.style={enhanced,breakable,boxrule=0.9pt,arc=2.5pt,
  left=3mm,right=3mm,top=2.3mm,bottom=2.3mm,fonttitle=\bfseries,coltitle=white}}
% --- boîtes TUTORIEL (noms EXACTS, ne pas renommer) ---
\newtcolorbox{cle}[1][]{boxbase,colback=defcol!6,colframe=defcol,
  title={\textsc{À retenir}\ifx\relax#1\relax\relax\else\textmd{ : \itshape #1}\fi}}
\newtcolorbox{methode}[1][]{boxbase,colback=methcol!7,colframe=methcol,sharp corners,
  title={\textsc{Méthode}\ifx\relax#1\relax\relax\else\textmd{ --- \itshape #1}\fi}}
\newtcolorbox{exemplebox}[1][]{boxbase,colback=excol!5,colframe=excol,
  title={\textsc{Exemple}\ifx\relax#1\relax\relax\else\textmd{ : \itshape #1}\fi}}
\newtcolorbox{piege}[1][]{boxbase,colback=attncol!6,colframe=attncol,
  title={\textsc{Piège}\ifx\relax#1\relax\relax\else\textmd{ : \itshape #1}\fi}}
\newtcolorbox{remarque}[1][]{boxbase,colback=black!4,colframe=black!45,
  title={\textsc{Remarque}\ifx\relax#1\relax\relax\else\textmd{ : \itshape #1}\fi}}
% --- boîtes EXERCICE / CORRIGÉ (noms EXACTS, auto-comptées) ---
\newtcolorbox[auto counter]{exo}[1][]{boxbase,colback=primary!4,colframe=primary,
  title={\textsc{Exercice}~\thetcbcounter\ifx\relax#1\relax\relax\else\textmd{ --- \itshape #1}\fi}}
\newtcolorbox[auto counter]{corrige}[1][]{boxbase,colback=excol!4,colframe=excol,
  title={\textsc{Corrigé}~\thetcbcounter\ifx\relax#1\relax\relax\else\textmd{ --- \itshape #1}\fi}}
\newcommand{\R}{\mathbb{R}}\newcommand{\N}{\mathbb{N}}\newcommand{\Z}{\mathbb{Z}}\newcommand{\ds}{\displaystyle}
% (un \usepackage{fancyhdr}+en-tête est possible ; il est ignoré côté web)
% =========================================================================
\begin{document}
\newcommand{\AX}[2]{\ensuremath{\mathrm{AX_{#1}E_{#2}}}}
\begin{exo}[l'eau \ce{H2O}]
L'oxygène central est lié à $2$ \ce{H} et porte $2$ doublets libres.
\begin{enumerate}[label=\alph*)]
  \item Écris la notation \AX{n}{m}.
  \item Donne la géométrie de répulsion \emph{et} la géométrie moléculaire.
  \item Quel est l'angle \ce{H-O-H} approximatif ?
\end{enumerate}
\end{exo}

\begin{exo}[l'ammoniac \ce{NH3}]
L'azote central est lié à $3$ \ce{H} et porte $1$ doublet libre.
\begin{enumerate}[label=\alph*)]
  \item Notation \AX{n}{m}.
  \item Géométrie de répulsion et géométrie moléculaire.
  \item Angle \ce{H-N-H} approximatif (plus grand ou plus petit que celui de \ce{H2O} ?).
\end{enumerate}
\end{exo}

\begin{exo}[le dioxyde de carbone \ce{CO2}]
Le carbone central est lié à $2$ oxygènes par deux doubles liaisons ; il n'a aucun doublet libre.
\begin{enumerate}[label=\alph*)]
  \item Combien de sites de répulsion ? Notation \AX{n}{m}.
  \item Géométrie de répulsion et géométrie moléculaire.
  \item Angle \ce{O-C-O}.
\end{enumerate}
\end{exo}

\begin{exo}[le dioxyde de soufre \ce{SO2}]
Le soufre central est lié à $2$ oxygènes et porte $1$ doublet libre.
\begin{enumerate}[label=\alph*)]
  \item Notation \AX{n}{m}.
  \item Géométrie de répulsion et géométrie moléculaire.
  \item L'angle \ce{O-S-O} est-il exactement $120^\circ$ ? Justifie.
\end{enumerate}
\end{exo}

\begin{exo}[lire une géométrie]
Voici la molécule de trifluorure de bore \ce{BF3}.
\begin{center}
\begin{tikzpicture}[scale=1.0]
  \definecolor{atomB}{RGB}{230,160,160}
  \definecolor{atomF}{RGB}{150,210,120}
  \coordinate (B) at (0,0);
  \coordinate (F1) at (90:1.5);
  \coordinate (F2) at (210:1.5);
  \coordinate (F3) at (330:1.5);
  \draw[line width=1.6pt,black!70] (B)--(F1);
  \draw[line width=1.6pt,black!70] (B)--(F2);
  \draw[line width=1.6pt,black!70] (B)--(F3);
  \draw[primarydark] (90:0.62) arc (90:210:0.62);
  \node[primarydark,font=\scriptsize] at (150:0.95) {$120^\circ$};
  \fill[atomB,draw=black!50] (B) circle (0.40); \node[font=\bfseries] at (B){B};
  \foreach \p in {F1,F2,F3}{\fill[atomF,draw=black!40] (\p) circle (0.34);
    \node[font=\footnotesize\bfseries] at (\p){F};}
\end{tikzpicture}
\end{center}
\begin{enumerate}[label=\alph*)]
  \item Combien de doublets libres porte le bore ? En déduis la notation \AX{n}{m}.
  \item Donne la géométrie de répulsion et la géométrie moléculaire.
  \item L'atome de bore respecte-t-il l'octet ?
\end{enumerate}
\end{exo}

\end{document}
